\documentclass[unicode,11pt,a4paper,oneside,numbers=endperiod,openany]{scrartcl}
\usepackage{amsmath}
\usepackage{amssymb}
\usepackage{ifthen}
\usepackage[utf8]{inputenc}
\usepackage{graphics}
\usepackage{graphicx}
\usepackage{hyperref}

\pagestyle{plain}
\voffset -5mm
\oddsidemargin  0mm
\evensidemargin -11mm
\marginparwidth 2cm
\marginparsep 0pt
\topmargin 0mm
\headheight 0pt
\headsep 0pt
\topskip 0pt        
\textheight 255mm
\textwidth 165mm

\newcommand{\duedate} {}
\newcommand{\setduedate}[1]{%
\renewcommand\duedate {Due date:~ #1}}
\newcommand\isassignment {false}
\newcommand{\setassignment}{\renewcommand\isassignment {true}}
\newcommand{\ifassignment}[1]{\ifthenelse{\boolean{\isassignment}}{#1}{}}
\newcommand{\ifnotassignment}[1]{\ifthenelse{\boolean{\isassignment}}{}{#1}}

\newcommand{\assignmentpolicy}{
\begin{table}[h]
\begin{center}
\scalebox{0.8} {%
\begin{tabular}{|p{0.02cm}p{16cm}|}
\hline
&\\
\multicolumn{2}{|c|}{\Large\textbf{HPC Lab for CSE 2024 ---  Submission Instructions}}\\
\multicolumn{2}{|c|}{\large\textbf{(Please, notice that following instructions are mandatory: }}\\
\multicolumn{2}{|c|}{\large\textbf{submissions that don't comply with, won't be considered)}}\\
&\\
\textbullet & Assignments must be submitted to \href{https://moodle-app2.let.ethz.ch/course/view.php?id=22516}{Moodle} (i.e. in electronic format).\\
\textbullet & Provide both executable package and sources (e.g. C/C++ files, Matlab). 
If you are using libraries, please add them in the file. Sources must be organized in directories called:\\
\multicolumn{2}{|c|}{\textit{Project\_number\_lastname\_firstname}}\\
& and  the  file must be called:\\
\multicolumn{2}{|c|}{\textit{project\_number\_lastname\_firstname.zip}}\\
\multicolumn{2}{|c|}{\textit{project\_number\_lastname\_firstname.pdf}}\\
\textbullet &  The TAs will grade your project by reviewing your project write-up, and looking at the implementation 
                 you attempted, and benchmarking your code's performance.\\

\textbullet & You are allowed to discuss all questions with anyone you like; however: (i) your submission must list anyone you discussed problems with and (ii) you must write up your submission independently.\\
\hline
\end{tabular}
}
\end{center}
\end{table}
}
\newcommand{\punkte}[1]{\hspace{1ex}\emph{\mdseries\hfill(#1~\ifcase#1{Points}\or{Points}\else{Points}\fi)}}


\newcommand\serieheader[6]{
\thispagestyle{empty}%
\begin{flushleft}
\includegraphics[width=0.4\textwidth]{ETHlogo_13}
\end{flushleft}
  \noindent%
  {\large\ignorespaces{\textbf{#1}}\hspace{\fill}\ignorespaces{ \textbf{#2}}}\\ \\%
  {\large\ignorespaces #3 \hspace{\fill}\ignorespaces #4}\\
  \noindent%
  \bigskip
  \hrule\par\bigskip\noindent%
  \bigskip {\ignorespaces {\Large{\textbf{#5}}}
  \hspace{\fill}\ignorespaces \large \ifthenelse{\boolean{\isassignment}}{\duedate}{#6}}
  \hrule\par\bigskip\noindent%  \linebreak
 }

\makeatletter
\def\enumerateMod{\ifnum \@enumdepth >3 \@toodeep\else
      \advance\@enumdepth \@ne
      \edef\@enumctr{enum\romannumeral\the\@enumdepth}\list
      {\csname label\@enumctr\endcsname}{\usecounter
        {\@enumctr}%%%? the following differs from "enumerate"
	\topsep0pt%
	\partopsep0pt%
	\itemsep0pt%
	\def\makelabel##1{\hss\llap{##1}}}\fi}
\let\endenumerateMod =\endlist
\makeatother




\usepackage{textcomp}






\begin{document}


\setassignment
\setduedate{Wednesday 26 June 2024, 23:59}

\serieheader{AI in the Sciences and Engineering}{Spring Semester 2024}
            {Student: Carla Judith L\'opez Zurita}
            {}{Project 3}{}
\newline

The main objective of the project is to apply the concepts learned in class
related to Neural Differential Equations and Hybrid Workflows.


\section{Inverted pendulum}\label{sec:task1}
The objective of this task is to control the inverted pendulum problem using a
Neural Network. The inverted pendulum is a classic problem in control theory,
where the goal is to balance a pendulum in an upright position by applying an 
external force $F$ to the cart that holds the pendulum. 
The system is described by the following differential equation:
\begin{equation}
    (M+m) \ddot{x} + m l \ddot{\theta} \cos(\theta) - m l \dot{\theta}^2 \sin(\theta) = F
    l \ddot{\theta} + g \sin(\theta) = 0
\end{equation}
where $\theta$ is the angle of the pendulum with respect to the vertical,
$g=9.81$ m/s$^2$ is
the acceleration due to gravity, $l=1.0$ m is the length of the pendulum,
$m=0.1$ kg is the
mass of the pendulum, $M=1.0$ kg is the mass of the cart and $F$ is the force
applied to the cart in Newtons.

\subsection*{Solving the coupled ODE system}
First, we need to solve the coupled ODE system to simulate the dynamics of the
inverted pendulum. I used Runge-Kutta 4th order method to solve the system.
We simulated the system with a sinusoidal force $F(t) = 10 \sin(t)$.
The results are shown in Figure~\ref{fig:ode}.
\begin{figure}[h]
    \centering
    \includegraphics[width=0.8\textwidth]{figures/ode.png}
    \caption{Simulation of the inverted pendulum with a sinusoidal force.}
    \label{fig:ode}
\end{figure}

\subsection*{Learning to balance the pendulum}

\section{Task 2}\label{sec:task2}


\section{Bonus Task}\label{sec:task3}


\bibliographystyle{plain}
\bibliography{bibliography}

\end{document}
